\documentclass[11pt,a4paper]{article}

\usepackage[UTF8]{ctex}
\usepackage{amsmath,amssymb}
\usepackage{geometry}
\geometry{margin=2.5cm}

\title{3 状态横向控制器（LQR）推导}
\author{}
\date{\today}

\begin{document}
\maketitle

\section{误差定义（VST 系投影）}

\begin{equation}
\begin{bmatrix}e_x\\ e_y\end{bmatrix}
= R_z(-\theta_{vst})\cdot\big(\vec p_{vst}-\vec p_{car}\big),
\qquad
R_z(\theta)=\begin{bmatrix}\cos\theta & -\sin\theta\\ \sin\theta & \cos\theta\end{bmatrix}
\end{equation}

\begin{align}
e_\theta &= \mathrm{wrap}_\pi\big(\theta_{vst}-\theta_{car}\big) \\
\dot e_\theta &= \frac{v}{L}\tan\delta_{vst} - \omega_{gyro}
\end{align}

\section{误差动力学（误差空间）}

\begin{align}
\dot e_y   &\approx v\,e_\theta \qquad \text{（小角度）} \\
\dot e_\theta &= \dot e_\theta \\
\ddot e_\theta &= -w_o\,\dot e_\theta - w_o\,\frac{v}{L}\,\Delta\delta
\end{align}

\section{状态空间}

\begin{equation}
x=\begin{bmatrix}e_y\\ e_\theta\\ \dot e_\theta\end{bmatrix},\quad u=\Delta\delta,\quad
\dot x =
\underbrace{\begin{bmatrix}
0 & v & 0\\
0 & 0 & 1\\
0 & 0 & -w_o
\end{bmatrix}}_{A}
x +
\underbrace{\begin{bmatrix}
0\\0\\-w_o\frac{v}{L}
\end{bmatrix}}_{B}
u
\end{equation}

\section{LQR}

\begin{samepage}

\begin{equation}
J=\int_0^\infty \left(x^T Q x + R u^2\right)\mathrm{d}t,
\qquad
Q=\mathrm{diag}(q_{ey},\,0,\,0),\quad R=r
\end{equation}

\begin{align}
u^* &= -K_{ric}\,x \\
K_{ric} &= R^{-1}B^T P \\
A^T P + P A - P B R^{-1} B^T P + Q &= 0
\end{align}

\end{samepage}

\section{增益表}

存储约定 $K=-K_{ric}$，故 $u=+Kx$。$K(v)=\big[K_{ey}(v),\,K_{e\theta}(v),\,K_{\dot e\theta}(v)\big]$
离线由 Riccati 方程按工作速度范围求解成表，车端二分 + 线性插值。

\section{控制律}

\begin{equation}
\Delta\delta_{fb}=K_{ey}\,e_y + K_{e\theta}\,e_\theta + K_{\dot e\theta}\,\dot e_\theta,
\qquad
\delta_{cmd}=\mathrm{clamp}\big(\delta_{vst}+\Delta\delta_{fb}\big)
\end{equation}

\end{document}
